\documentclass[conference]{IEEEtran}
\usepackage{cite}
\usepackage{amsmath,amssymb}
\usepackage{graphicx}
\usepackage{booktabs}
\usepackage{tikz}
\usetikzlibrary{arrows.meta,positioning}
\usepackage{url}
\usepackage[hidelinks]{hyperref}
\graphicspath{{figures/}}
\hypersetup{
  pdftitle={Automated Memory Forensics for Malware Detection and Incident Response},
  pdfauthor={Joel C. Okore, Andrei Petrovski, Igor V. Kotenko},
  pdfkeywords={memory forensics, incident response, security automation, malware detection}
}

\begin{document}

\title{Automated Memory Forensics for\\
Malware Detection and Incident Response}

\author{\IEEEauthorblockN{Joel C. Okore}
\IEEEauthorblockA{\textit{Faculty of Computer Science} \\
\textit{Innopolis University}\\
Innopolis, Tatarstan, Russia \\
j.okore@innopolis.university}
\and
\IEEEauthorblockN{Andrei Petrovski}
\IEEEauthorblockA{\textit{Faculty of Computer Science} \\
\textit{Innopolis University}\\
Innopolis, Tatarstan, Russia \\
a.petrovski@innopolis.ru}
\and
\IEEEauthorblockN{Igor V. Kotenko}
\IEEEauthorblockA{\textit{Security Problems Lab} \\
\textit{SPC RAS}\\
St. Petersburg, Russia \\
ivkote@comsec.spb.ru}
\IEEEauthorblockA{\textit{Institute of Information Security} \\
\textit{Innopolis University}\\
Innopolis, Tatarstan, Russia \\
i.kotenko@innopolis.ru}}

\maketitle

\begin{abstract}
Memory forensics can provide evidence of malicious activity that is unavailable from endpoint alerts alone, but acquisition, analysis, and response often require separate operational steps. This paper presents an automated memory-forensics pipeline that connects these steps through an event-driven workflow. Wazuh alerts trigger remote memory acquisition with WinPMEM; file-integrity events initiate Volatility3 for memory analysis and extraction of forensic artifacts, which will then be used as features for machine learning classification; and a malicious prediction invokes network isolation and an analyst notification. The implementation uses a Windows endpoint, an Ubuntu manager, and explicit file and message interfaces between the stages of the pipeline. The case study in consideration demonstrates the workflow from an endpoint alert through memory acquisition, classification, and automated response. In the isolated experiment, WannaCry ransomware was detonated on a monitored endpoint, the pipeline was triggered automatically and progressed from memory acquisition to response initiation in under three minutes. Execution logs and a subsequent connectivity test document the observed response behavior. Supporting experiments examine classifier portability and a host-relative scoring extension, identifying requirements for calibration and clean-reference integrity. The main contribution is the design, implementation and integration of the forensic response workflow, together with the operational lessons needed to extend it.
\end{abstract}

\begin{IEEEkeywords}
memory forensics, incident response, security automation, malware detection, Wazuh, Volatility3
\end{IEEEkeywords}

\section{Introduction}

Security operations depend on both detecting suspicious activity and obtaining the evidence needed to investigate it. Endpoint logs can indicate an unusual command or process, while a memory image provides a basis for examining the system state associated with that activity. Because volatile artifacts can disappear when processes terminate or memory is reused, initiating acquisition is an operational concern.

A practical memory forensics workflow must coordinate several tools. An alert must identify the endpoint, a privileged acquisition utility must obtain its memory image, the image must become accessible to an analysis host, and the analysis result must reach a response mechanism. When these transitions require manual intervention, an analyst must repeatedly transfer context between detection, forensic analysis, and containment. Automating the transitions provides a way to connect forensic collection to an existing incident response workflow.

In this work, we design and implement such a workflow using Wazuh active response~\cite{wazuh}, WinPMEM~\cite{winpmem}, Volatility3~\cite{volatility}, and a machine learning classifier. Wazuh coordinates the workflow: endpoint alerts trigger memory acquisition, while file-integrity events initiate feature extraction from the captured image and classification of the resulting feature vector. The classifier returns a structured decision that the response script uses to request network isolation and notify the security team. The system therefore connects collection, analysis, and response action using interfaces available in an open-source security stack.

The application setting is a Windows endpoint monitored by an Ubuntu Wazuh manager. We evaluate functional integration through an isolated WannaCry case study, examining the artifacts and logs produced as the workflow executes. A supporting assessment of the classifier and a host-relative extension addresses what additional evidence is needed before transferring automated response to different hosts and workloads.

The contributions are:
\begin{enumerate}
\item An event-driven architecture that connects endpoint detection to remote memory acquisition, forensic feature extraction, and automated response. 
\item A working implementation with explicit artifact and message interfaces, demonstrated through a WannaCry ransomware case study with execution logs and observed connectivity loss after endpoint isolation action.
\item An assessment of deployment requirements, including stage completion, response verification, classifier portability, and baseline integrity.
\end{enumerate}

\section{Related Work}

Wazuh active response associates scripts with detection rules, severity levels, or rule groups~\cite{wazuh}. WinPMEM supports acquisition of physical memory into raw images~\cite{winpmem}, while Volatility3 provides a framework for extracting forensic artifacts from memory~\cite{volatility}. These tools supply the monitoring, acquisition, and analysis capabilities used here. Our work concerns their integration into a coordinated workflow whose analysis output can drive a response action.

Memory-feature classification offers one way to automate the interpretation stage. VolMemLyzer~\cite{lashkari2021} extracts features for malware classification, and CIC-MalMem-2022~\cite{carrier2022} extends this approach with obfuscation-related memory features. Carrier et al. report 99.00\% accuracy and a 99.02\% F1-score for a stacked ensemble. Such results motivate the use of classifiers within forensic workflows, while leaving their behavior on a deployment endpoint to be evaluated.

Arp et al.~\cite{arp2022} discuss sampling bias, data leakage, and evaluation pitfalls in security machine learning. We consider these issues as part of assessing the decision component of an operational system. The central question in this paper is how the forensic response workflow is implemented and behaves in the testbed; benchmark analysis supports the discussion of its transfer to other environments.

\section{System Design and Implementation}
\label{sec:architecture}

\subsection{Operational Objectives}
The system is designed around three objectives: initiate memory dump collection triggered by an endpoint alert, pass the resulting evidence through analysis without manual invocation, and connect the resulting decision of the analysis to endpoint containment and analyst notification. Memory acquisition precedes network isolation so that the manager can obtain and analyze the image while its management and data paths remain available. 
% This ordering makes acquisition part of the response workflow and also introduces a delay before containment.

The Windows 10 endpoint runs the Wazuh agent and WinPMEM. It is configured to supply PowerShell operational events and exposes the memory dump directory through an SMB-shared network folder. The Ubuntu manager hosts Wazuh rules, active-response scripts, Volatility3, the serialized classifier, and the notification client. Memory dump acquisition and endpoint isolation use Windows Remote Management (WinRM), invoked through NetExec. The image is made available to the manager through the SMB-shared network directory. Fig.~\ref{fig:architecture} shows the component placement and system architecture.

The implementation source code is available in the project repository~\cite{projectrepo}.

\begin{figure}[!t]
\centering
\includegraphics[width=0.50\textwidth]{arch_diag.png}
\caption{Deployment architecture and script-level workflow across the Windows
endpoint and Ubuntu manager.
% AR denotes Wazuh active response; SOAR denotes the automated response stage. Network quarantine is implemented by disabling endpoint network adapters over WinRM.
}
\label{fig:architecture}
\end{figure}

\subsection{Alert-Driven Memory Acquisition}
The Wazuh agent collects the Windows PowerShell operational log. Custom rules were configured to match indicators of compromise such as encoded-command execution and known offensive cmdlets. Wazuh active response scripts associated with these rules trigger memory dump acquisition on the monitored host and are configured with a severity level of 10. This stage uses the endpoint alert to trigger a memory dump of the monitored host, and subsequent classification is based on features extracted from the captured memory image.

The acquisition script reads the Wazuh JSON alert from standard input and extracts the agent name and IP address. It then invokes NetExec over WinRM to launch the preinstalled WinPMEM executable on that endpoint. The raw image is written to the configured dump path, and command output is appended to the acquisition log. Reading the target address from the alert connects the collection request to the endpoint that produced the triggering event.

The shared directory links the Windows acquisition path to the manager's analysis path. 
% Its use allows the analysis utilities to run on the manager while WinPMEM operates on the endpoint. 
The prototype uses fixed image and feature filenames suitable for a single-endpoint experiment. Concurrent incidents from different endpoints would require separate artifact locations and identifiers.

\subsection{Event-Driven Orchestration}
Wazuh File Integrity Monitoring (FIM) coordinates the pipeline by detecting changes to the memory image and feature CSV. Changes to the raw memory image trigger feature extraction through rules 100100/100101. Changes to the feature CSV invoke the response script through rules 100102/100103; this script classifies the extracted features and applies the response policy. These custom rules are defined specifically for the pipeline and are not included in Wazuh's default ruleset. Their identifiers were chosen arbitrarily, and the active-response configuration associates each rule with the corresponding pipeline script. The file-change events determine when each stage runs, while the files provide the data it processes. Table~\ref{tab:interfaces} summarizes the inputs and
outputs of each component.

\begin{table}[t]
\caption{Pipeline Components and Data Interfaces}
\label{tab:interfaces}
\centering
\footnotesize
\begin{tabular}{@{}p{1.45cm}p{2.85cm}p{3.55cm}@{}}
\toprule
Component & Input or trigger & Output or action \\
\midrule
Acquisition & Wazuh alert JSON & Remote WinPMEM invocation; raw memory image \\
Extraction & Image-file FIM event & Forensic feature vector in CSV format \\
Classifier & Feature vector & Class label and model score in JSON \\
Response & Feature-file FIM event and prediction & Network isolation request; analyst notification \\
\bottomrule
\end{tabular}
\end{table}

% The interfaces define three representations of an incident: the captured memory state, the trigger event, and the decision supplied to the response policy. Retaining these intermediate artifacts supports inspection of the relationship between evidence and action. It also permits changes to feature extraction or classification without redesigning acquisition, provided that the associated input and output interfaces remain compatible.

% The current implementation uses file modifications as transition events. This provides a direct integration with Wazuh FIM, but does not enforce file completion or unique event processing. The corresponding requirements for completion detection and duplicate suppression are addressed in Section~\ref{sec:discussion}.

The pipeline produces a memory image, a CSV file of extracted features, and a classification result. Together with the recorded Wazuh events, these artifacts allow an analyst to trace how captured evidence leads to a response decision. Defined file and message interfaces also allow the feature extractor or classifier to be updated without changing memory acquisition, provided that each component continues to accept the expected inputs and produce the outputs required by the next stage.

In the current implementation, a file modification triggers the next stage without first verifying that writing has finished. Feature extraction or classification may therefore begin before its input file is complete. Repeated events for the same file may also invoke a stage more than once. Section~\ref{sec:discussion} discusses the completion checks and duplicate-event
handling needed to address these limitations.


\subsection{Forensic Feature Representation and Classification}
Classification uses two features derived from Volatility3 service-scan output: the number of service records, \texttt{svcscan.nservices}, and the number of distinct nonzero process identifiers associated with those records, \texttt{svcscan.process\_services}. This aggregation converts the plugin's variable-length output into a fixed-length representation of Windows service state and confines extraction to a single plugin.

An XGBoost model performs binary classification of the resulting vector. Preprocessing uses a standardization transform fitted on the training partition and serialized with the model. Inference applies the same transform and feature ordering. The original two-feature model achieved 99.84\% accuracy on a 70/30 split of the pre-preocessed CIC-MalMem-2022 dataset. The implications of this benchmark result for deployment are examined in Section~\ref{sec:assessment}.

The decision interface consists of a JSON message containing the numeric label, class name, and malicious-class model score. The predicted label controls the response policy; the score provides additional context in the analyst notification. The implementation initiates containment for label one and leaves the endpoint unchanged for a benign prediction. A JSON parsing failure terminates this invocation without a containment request.

\subsection{Response Policy and Operational Records}
The response policy combines endpoint containment with analyst notification. A malicious classification initiates a WinRM command that disables the endpoint's network adapters. Notification is handled on the Ubuntu manager, which sends the affected host, predicted class, and classification score to Slack after requesting isolation. This placement allows notification to proceed independently of the endpoint's network availability following containment.

WinRM credentials and the notification webhook are supplied through environment variables. Operational records cover acquisition, extraction, classification, and response, providing evidence of execution at each stage. The distinction between command execution and its effect is evaluated through both response logs and endpoint connectivity observations in Section~\ref{sec:case}.

\section{Experimental Evaluation}
\label{sec:case}

\subsection{Testbed and Evaluation Criteria}
The evaluation uses an isolated two-VM testbed comprising a Windows 10 endpoint and an Ubuntu Wazuh manager. PowerShell event collection and FIM supply the detection inputs; WinRM and an SMB share provide the control and artifact-access paths. In the evaluation scenario, WannaCry ransomware was introduced through a PowerShell download-and-execution sequence which triggered automated memory dump acquisition on the monitored endpoint.
% from a server within the testbed.

Functional evaluation considers three properties: automatic invocation of successive stages, propagation of the acquired evidence through classification, and execution of the configured response. Wazuh events and script outputs provide evidence of stage invocation and artifact processing. Response logs and a subsequent endpoint connectivity test assess the containment and notification behavior. Workflow latency is measured between subsequent stages of the pipeline.
% acquisition-script and response-script start times.
The experiment evaluates integration in one malware scenario; the scope of generalization is discussed in Section~\ref{sec:discussion}.

\subsection{Functional Integration and Response}
The WannaCry experiment demonstrates automatic progression through acquisition, feature extraction, classification, and response. Raw-image and feature-file FIM events correspond to the configured extraction and response rules, while successful Volatility3 dump analysis and artifact extraction establishes that the acquired image reached the analysis stage. The resulting feature vector produced a malicious classification and activated the containment policy.

Table~\ref{tab:observations} relates these results to the evaluation criteria. The response output in Fig.~\ref{fig:response} records an isolation request and successful Slack notification. In the endpoint connectivity test, all four outbound probes from the isolated endpoint failed with local transmission errors, confirming that the network adapters were disabled as a result of the isolation.
% These observations support the functional operation of the response path in the testbed, without establishing coverage of every possible communication path.

\begin{figure*}[t]
\centering
\includegraphics[width=\textwidth]{response_execution.png}
\caption{Prediction and response output from the WannaCry experiment. The classifier's prediction of the memory image as malicious activates endpoint isolation and analyst notification. The displayed
probability is the model output for this capture.}
\label{fig:response}
\end{figure*}

\begin{table}[t]
\caption{Functional Evaluation of the Integrated Pipeline}
\label{tab:observations}
\centering
\footnotesize
\begin{tabular}{@{}p{2.5cm}p{5.55cm}@{}}
\toprule
Evaluation criterion & Experimental evidence \\
\midrule
Automatic stage invocation & Acquisition trigger and image/feature FIM events activate the configured stages \\
Evidence processing & Successful service scan, feature extraction, and classification \\
Containment & Isolation request and four failed endpoint connectivity probes \\
Analyst notification & Successful Slack delivery reported by the response client \\
\bottomrule
\end{tabular}
\end{table}

The separation of the classifier from response policy makes the feature vector and predicted label available for inspection at the decision boundary.
% The current decision message uses JSON; the experiment's retained screenshot uses an earlier dictionary-style serialization of the decision fields. This format difference does not alter the label-based response policy.

\subsection{Workflow Latency}
The pipeline reached the response stage 171\,s (2\,min 51\,s) after the acquisition stage started. Table~\ref{tab:timing} reports the available timestamps for the start of each stage. The acquisition-to-extraction interval is 60\,s, followed by 111\,s between extraction and response initiation. These intervals include both stage processing and the associated orchestration overhead.

\begin{table}[t]
\caption{Stage-Start Timing in the WannaCry Experiment}
\label{tab:timing}
\centering
\small
\begin{tabular}{@{}lrr@{}}
\toprule
Stage & Start time & Elapsed from acquisition \\
\midrule
Acquisition & 23:44:09 & 0\,s \\
Feature extraction & 23:45:09 & 60\,s \\
Prediction and response & 23:47:00 & 171\,s \\
\bottomrule
\end{tabular}
\par\smallskip
\begin{minipage}{\columnwidth}\footnotesize
% Timestamps: 16 May 2025, MSK. Elapsed time is measured between stage starts.
\end{minipage}
\end{table}

The measured interval spans the start of memory acquisition to the start of the response stage.
% Initial malicious execution and completed containment lack separate timestamps, preventing a full detection-to-containment latency estimate.
The experiment demonstrates automated execution of the integrated workflow in the testbed. Repeated trials with different memory sizes, workloads, and numbers of concurrently processed endpoints are
needed to assess how response times vary across operating conditions.

\section{Classifier Assessment}
\label{sec:assessment}

\subsection{Benchmark Results and Portability}
% Classifier portability is relevant to the system because its prediction determines endpoint isolation. A supporting evaluation therefore examines whether the benchmark result establishes a suitable basis for that decision across different hosts. CIC-MalMem-2022 contains 58{,}596 balanced records and 55 numeric features~\cite{carrier2022,cicmalmem}. The ablation uses 52 nonconstant predictors, a stratified 70/30 split with seed 7, and XGBoost models with 200 estimators and maximum depth 6. Standardization is fitted on the training partition.

% Accuracy reaches 99.99\% with all nonconstant features, 99.98\% with six host-configuration counters, and 99.84\% with the two service features. A descriptive rule classifying service counts at most 391 or above 1000 as malicious achieves 99.63\% on the full corpus. The rule is a corpus-level comparison rather than a held-out estimate. The limited difference between feature groups indicates that benchmark discrimination can be achieved using host-state counters alone.

% Distributional diagnostics further constrain the interpretation of these results. All 22 selected count and count-derived indicators have higher benign medians, and one-to-one load matching retains 148 records. Five-fold logistic regression places only 0.92\% of records in the intermediate probability band $0.1<P(\mathrm{malicious}\mid x)<0.9$. This statistic describes polarized class predictions, without isolating the effect of collection conditions. The reported use of SMOTE during dataset construction introduces an additional risk of dependencies between random partitions~\cite{carrier2022}. These findings motivate evaluation on representative endpoint captures before benchmark-trained decisions are used for operational containment.

Because a malicious prediction triggers endpoint isolation, the classifier must be evaluated on hosts beyond those represented in its training data. The supporting analysis compares feature sets to determine whether high benchmark accuracy can be reproduced using a small set of host-state measurements. The dataset used in this study, CIC-MalMem-2022 contains 58{,}596 records, equally divided between benign and malicious classes, and 55 numeric features~\cite{carrier2022,cicmalmem}. Removing constant features leaves 52 predictors. The trained XGBoost models use 200 estimators, a maximum tree depth of 6, and a stratified 70/30 training/test split with random seed 7. Standardization parameters are fitted using only the training data. 

Test accuracy is 99.99\% with all 52 predictors, 99.98\% with six host-configuration counters, and 99.84\% with the two service features. A simple rule based on service count alone labels a record as malicious when the count is 391 or lower, or greater than 1000. This rule correctly classifies 99.63\% of records in the entire dataset.
% ; this percentage is a descriptive result, not an independent test estimate
The similar model accuracies show that a small set of host-state counters can separate the dataset's classes almost as well as the full feature set. These results alone do not establish whether the classifier detects malicious activity across different host configurations.

Distributional diagnostics further constrain the interpretation of these results. All 22 selected count and count-derived indicators have higher benign medians, and one-to-one load matching retains 148 records. Five-fold logistic regression places only 0.92\% of records in the intermediate probability band $0.1<P(\mathrm{malicious}\mid x)<0.9$. This statistic describes polarized class predictions, without isolating the effect of collection conditions. The reported use of SMOTE during dataset construction introduces an additional risk of dependencies between random partitions~\cite{carrier2022}. These findings motivate evaluation on representative endpoint captures before benchmark-trained decisions are used for operational containment.


\subsection{Host-Relative Feature Extension}
A baseline-differential representation is evaluated as a possible extension to the classifier interface. It expresses a capture relative to clean reference observations from the same host, using 44 aggregates from eight Volatility3 plugins. Let $m_i$ denote the reference median of feature $i$ and $d_i$ its median absolute deviation multiplied by 1.4826. The feature deviation and aggregate score are
\begin{equation}
z_i(x)=\frac{x_i-m_i}{s_i+10^{-9}},\qquad
S_F(x)=\max_{i\in F}|z_i(x)|,
\end{equation}
where $F$ is the selected feature set. The scale is $s_i=d_i$ unless $d_i<10^{-6}$, in which case $s_i=\max(0.01|m_i|,1)$. This fallback provides a finite scale for features with negligible reference dispersion. The extension is assessed separately from the deployed XGBoost response path.

The pilot comprises ten sequential pre/post pairs from one Windows Server 2022 VM using Atomic Red Team~\cite{atomicredteam} tests associated with MITRE ATT\&CK T1055.001, .002, .003, .004, and .012~\cite{mitre1055}. The first six pairs form the retained evaluation set. Pre-test values of \texttt{malfind.commitCharge} shift from 259--265 in trials 1--6 to 773--778 in trials 7--10, indicating a persistent change in the reference state. This is consistent with incomplete cleanup, although the aggregate data do not identify the responsible process.

For the restricted set of 11 \texttt{malfind} and \texttt{ldrmodules} features, scores increase from 0.797 to 3.000 in the APC-injection pair and from 4.000 to 5.000 in the process-hollowing pair. No retained post-test capture reaches the default threshold of 6. The full 44-feature score flags three post-test captures and one reference capture. Because the six pre-test captures also fit the reference model, these observations characterize score behavior rather than independent detection and false-positive rates. The application requirement is therefore both representative baseline coverage and independent calibration of the response threshold.

\section{Discussion}
\label{sec:discussion}

\subsection{Operational Value and Response Ordering}
The operational value of the system is that an endpoint alert can initiate forensic collection and carry the resulting evidence through to a response decision. The WannaCry case study demonstrates this progression without manual stage invocation. The acquired image, feature record, and classification output also support analyst review of the decision. Assessing the effect on analyst workload requires comparison with a manual workflow.

The order of these stages also affects how long a contaminated endpoint remains connected. The evaluated configuration collects and analyzes memory before isolation to preserve WinRM and SMB access. The contaminated endpoint therefore remains connected during this interval. The appropriate ordering depends on the incident policy and the availability of management paths that remain accessible after isolation.

\subsection{Reliability Requirements}
% Reliable analysis requires completed input files and protection against repeated invocation.
The prototype responds to file modifications, so a stage may start before its input is complete or run more than once. Completion signals or atomic publication of completed files would address the first issue; duplicate-event suppression would address the second. Incident-specific artifact paths are also needed to keep concurrent captures separate.

Containment reliability depends on verifying the command's effect on the endpoint. Currently, the isolation wrapper implementation does not check the remote command's exit status. An operational implementation should check the command exit status and independently verify that the intended network isolation took effect. The case-study connectivity observation supports the demonstrated response, but equivalent verification needs to be part of each response invocation.

\subsection{Evaluation Scope and Further Validation}
These reliability requirements must be evaluated beyond the current demonstration. The system evaluation covers one Windows endpoint and one WannaCry experiment, establishing functional integration in that setting. It does not determine throughput under concurrent load, latency variability, recovery from failed stages, or detection coverage across malware families.

The supporting host-relative pilot informs baseline design but does not validate a replacement classifier. Workload and technique are coupled in a fixed rotation, and retained evidence lacks raw memory images and per-test execution transcripts. Broader validation should therefore combine repeated system trials under varied workloads, independent classifier calibration using representative endpoint captures, and tests of recovery after acquisition or response failures.


\section{Conclusion}

This work demonstrates the integration of memory forensics into an automated incident-response workflow. The implemented system uses Wazuh active-response rules to coordinate memory acquisition, forensic analysis, classification, and containment across a Windows endpoint and an Ubuntu manager. In the isolated WannaCry experiment, the pipeline executed without manual stage invocation and reached response initiation in under three minutes from the start of memory acquisition. Execution logs document the isolation request and analyst notification, while subsequent connectivity observations support the effect of the containment action.

The system's contribution is to connect forensic evidence collection directly to endpoint response while making the memory image, extracted features, and classification result available for analyst inspection. The supporting
classifier assessment highlights the need to validate decisions on captures representative of the deployment environment. Further development will strengthen evaluation of performance and failure recovery across multiple endpoints and workloads.
% Balance the final-page columns; revisit after changes to manuscript length.
% \IEEEtriggeratref{5}
\newpage
\begin{thebibliography}{10}
\bibitem{wazuh} Wazuh, ``Active Response,'' \textit{Wazuh Documentation}.
Accessed: Sep. 9, 2026. [Online]. Available:
\url{https://documentation.wazuh.com/current/user-manual/capabilities/active-response/index.html}

\bibitem{winpmem} Velocidex, ``WinPmem: A Physical Memory Acquisition Tool.''
Accessed: Sep. 9, 2026. [Online]. Available:
\url{https://github.com/Velocidex/WinPmem}

\bibitem{volatility} Volatility Foundation, ``Volatility 3 Documentation.''
Accessed: Sep. 9, 2026. [Online]. Available:
\url{https://volatility3.readthedocs.io/en/latest/}

\bibitem{lashkari2021} A. H. Lashkari, B. Li, T. L. Carrier, and G. Kaur,
``VolMemLyzer: Volatile Memory Analyzer for Malware Classification using
Feature Engineering,'' in \textit{Proc. Reconciling Data Analytics,
Automation, Privacy, and Security: A Big Data Challenge (RDAAPS)},
2021, pp. 1--8, doi: 10.1109/RDAAPS48126.2021.9452028.

\bibitem{carrier2022} T. Carrier, P. Victor, A. Tekeoglu, and A. H. Lashkari,
``Detecting Obfuscated Malware using Memory Feature Engineering,'' in
\textit{Proc. 8th Int. Conf. Information Systems Security and Privacy
(ICISSP)}, 2022, pp. 177--188, doi: 10.5220/0010908200003120.

\bibitem{arp2022} D. Arp, E. Quiring, F. Pendlebury, A. Warnecke,
F. Pierazzi, C. Wressnegger, L. Cavallaro, and K. Rieck,
``Dos and Don'ts of Machine Learning in Computer Security,'' in
\textit{Proc. 31st USENIX Security Symp.}, 2022, pp. 3971--3988.

\bibitem{projectrepo} J. C. Okore, ``CCF-Project,'' GitHub repository.
[Online]. Available:
\url{https://github.com/Joellots/CCF-project}

\bibitem{cicmalmem} Canadian Institute for Cybersecurity,
``Malware Memory Analysis (CIC-MalMem-2022),'' University of New Brunswick.
Accessed: Sep. 9, 2026. [Online]. Available:
\url{https://www.unb.ca/cic/datasets/malmem-2022.html}

\bibitem{atomicredteam} Red Canary, ``Atomic Red Team.''
Accessed: Sep. 9, 2026. [Online]. Available:
\url{https://github.com/redcanaryco/atomic-red-team}

\bibitem{mitre1055} MITRE, ``Process Injection, Technique T1055,''
\textit{MITRE ATT\&CK}.
Accessed: Sep. 9, 2026. [Online]. Available:
\url{https://attack.mitre.org/techniques/T1055/}

\end{thebibliography}
\end{document}
